% =====================================================================
\chapter{Entropía, física y ajuste por preferencias}
\label{tema:5}
% =====================================================================

Este tema relaciona la entropía de Shannon con la entropía de la física estadística y muestra que la misma herramienta, la desigualdad de Gibbs, resuelve tres problemas en apariencia distintos: la distribución de máxima entropía con una restricción de energía media, que es la distribución de Boltzmann y, en particular, la función softmax; el principio variacional de la energía libre; y el problema de maximizar una recompensa con una penalización de Kullback--Leibler, que está en la base del ajuste de los LLM por preferencias humanas. Se sigue~\cite[cap.~12]{ct} y el artículo de Jaynes~\cite{jaynes} para la máxima entropía, el artículo de Landauer~\cite{landauer} para el coste físico del borrado y los artículos de Ouyang \emph{et al.}~\cite{instructgpt} y Rafailov \emph{et al.}~\cite{dpo} para el ajuste por preferencias. El tema «Applications of information theory in other sciences» de~\cite{camit} trata el mismo puente con la física.

En este tema los logaritmos son naturales y las entropías y divergencias se miden en nats, salvo indicación expresa.

\section{La entropía en la física estadística}

Un sistema físico macroscópico, descrito por unas pocas magnitudes (energía, volumen, número de partículas), es compatible con un número enorme de configuraciones microscópicas o \emph{microestados}.

\begin{definicion}[Entropías de Boltzmann y de Gibbs]\label{def:boltzmann}
Si un macroestado es compatible con $W$ microestados equiprobables, su \emph{entropía de Boltzmann} es $S=k_B\ln W$. Si los microestados tienen probabilidades $p_i$, la \emph{entropía de Gibbs} es
\[
S=-k_B\sum_ip_i\ln p_i,
\]
donde $k_B=1{,}380649\cdot10^{-23}\,\mathrm{J/K}$ es la constante de Boltzmann.
\end{definicion}

La entropía de Gibbs es $k_B$ veces la entropía de Shannon en nats, y la de Boltzmann es el caso uniforme, en el que $H=\ln W$. La entropía termodinámica de un macroestado es, salvo la constante $k_B$, la cantidad de información que falta para identificar el microestado~\cite{jaynes}.

\begin{center}
\begin{tabular}{@{}ll@{}}
\toprule
Física estadística & Teoría de la información \\
\midrule
microestado & mensaje o secuencia concreta \\
$S=k_B\ln W$ & $H=\log W$ bits (caso uniforme) \\
$S=-k_B\sum_ip_i\ln p_i$ & $H=-\sum_ip_i\log p_i$ \\
distribución de Boltzmann $e^{-E/k_BT}/Z$ & softmax de los logits con temperatura \\
energía libre $F=\langle E\rangle-TS$ & objetivo con penalización de Kullback--Leibler \\
\bottomrule
\end{tabular}
\end{center}

\begin{proposicion}\label{prop:binomial}
Para $0\le k\le n$, $\binom{n}{k}\le 2^{nH_b(k/n)}$, donde $H_b$ es la entropía binaria en bits.
\end{proposicion}

\addtocounter{definicion}{1}%

\section{Máxima entropía y distribución de Boltzmann}

Cuando de una distribución solo se conoce el valor medio de una magnitud, el principio de máxima entropía de Jaynes~\cite{jaynes} propone elegir, entre todas las distribuciones compatibles, la de mayor entropía: la que no añade ninguna información que no se tenga.

\begin{teorema}[Distribución de máxima entropía]\label{teo:maxent}
Sea $E\colon\X\to\R$ una función de energía sobre un alfabeto finito, sea $\lambda\in\R$ y sea
\[
p_\lambda(x)=\frac{e^{-\lambda E(x)}}{Z(\lambda)},\qquad Z(\lambda)=\sum_{x\in\X}e^{-\lambda E(x)}.
\]
Para toda distribución $p$ sobre $\X$ con la misma energía media, $\E_p[E]=\E_{p_\lambda}[E]$, se cumple $H(p)\le H(p_\lambda)$, con igualdad si y solo si $p=p_\lambda$.
\end{teorema}

En física, $\lambda=1/(k_BT)$ y $p_\lambda$ es la \emph{distribución de Boltzmann} a temperatura $T$: los estados de energía baja son más probables, y tanto más cuanto menor es la temperatura.

\addtocounter{definicion}{1}%

\section{Energía libre y principio variacional}

El problema siguiente aparece en física, con la energía libre, y en el ajuste de modelos de lenguaje, con una recompensa. Sean $\pref$ una distribución de referencia sobre un conjunto finito de respuestas $\mathcal{Y}$, con $\pref(y)>0$ para todo $y$, una recompensa $r\colon\mathcal{Y}\to\R$ y un coeficiente $\beta>0$. Se busca la distribución $\pi$ que maximiza
\[
\mathcal{F}(\pi)=\E_{y\sim\pi}[r(y)]-\beta\,\KL{\pi}{\pref}.
\]

\begin{teorema}[Principio variacional de Gibbs]\label{teo:variacional}
Sean $Z=\sum_y\pref(y)\,e^{r(y)/\beta}$ y $\pi^{*}(y)=\pref(y)\,e^{r(y)/\beta}/Z$. Para toda distribución $\pi$ sobre $\mathcal{Y}$,
\[
\mathcal{F}(\pi)=\beta\ln Z-\beta\,\KL{\pi}{\pi^{*}}.
\]
En consecuencia, $\pi^{*}$ es el único maximizador de $\mathcal{F}$ y el valor máximo es $\beta\ln Z$.
\end{teorema}

\begin{corolario}\label{cor:energia-libre}
Si $\pref$ es uniforme sobre $K$ respuestas, $\mathcal{F}(\pi)=\E_\pi[r]+\beta H(\pi)-\beta\ln K$. Con energía $E=-r$ y temperatura $\beta$, maximizar $\mathcal{F}$ equivale a minimizar la \emph{energía libre} $\E_\pi[E]-\beta H(\pi)$, cuyo minimizador es la distribución de Boltzmann.
\end{corolario}

\section{El ajuste por preferencias con penalización de Kullback--Leibler}

Tras el preentrenamiento y el ajuste supervisado, un modelo de lenguaje se ajusta a las preferencias humanas con aprendizaje por refuerzo (RLHF, módulo~21). Un \emph{modelo de recompensa} $r(x,y)$, entrenado con comparaciones humanas, puntúa la respuesta $y$ al mensaje $x$. Maximizar solo $r$ llevaría el modelo hacia respuestas extrañas que explotan los defectos del modelo de recompensa (\emph{reward hacking}). Por ello, Ouyang \emph{et al.}~\cite{instructgpt} añaden a la recompensa una penalización por alejarse del modelo supervisado $\pref$, que en esperanza es $\beta\,\KL{\pi}{\pref}$; para cada mensaje $x$ se maximiza la cantidad $\mathcal{F}$ del apartado anterior.

Por el \cref{teo:variacional}, la política óptima es $\pi^{*}(y\mid x)\propto\pref(y\mid x)\,e^{r(x,y)/\beta}$: la referencia «recalentada» por la recompensa, con $\beta$ en el papel de temperatura. En la práctica el conjunto de respuestas es inmenso y $\pi^{*}$ no se calcula de forma explícita, sino que se aproxima con métodos de gradiente; el teorema describe a qué apunta esa optimización.

\addtocounter{definicion}{1}%

\begin{observacion}[Optimización directa de preferencias]
Despejando la recompensa en la solución óptima, $r(x,y)=\beta\ln\bigl(\pi^{*}(y\mid x)/\pref(y\mid x)\bigr)+\beta\ln Z(x)$: cada política define implícitamente una recompensa. Rafailov \emph{et al.}~\cite{dpo} observan que, en el modelo de Bradley--Terry de comparaciones por pares, el término $\beta\ln Z(x)$ se cancela, y obtienen un método, DPO, que ajusta la política directamente con las comparaciones humanas, sin entrenar un modelo de recompensa separado (\cref{eje:5-dpo}).
\end{observacion}

\section{El coste físico de la información}

\begin{proposicion}[Principio de Landauer]\label{prop:landauer}
Borrar un bit de información en un sistema en contacto con un entorno a temperatura $T$ disipa en el entorno un calor de al menos $k_BT\ln 2$.
\end{proposicion}

La idea, debida a Landauer~\cite{landauer}, es la siguiente. Un bit de memoria cuyo valor se desconoce está en uno de dos estados equiprobables, con entropía $k_B\ln 2$; tras el borrado, que lo pone siempre a~0, su entropía es nula. Como la entropía total no puede disminuir, al menos $k_B\ln 2$ de entropía pasa al entorno, lo que exige un calor $Q\ge T\,\Delta S=k_BT\ln 2$. A $T=300$~K, este mínimo vale $1{,}380649\cdot10^{-23}\cdot300\cdot0{,}6931=2{,}87\cdot10^{-21}$~J por bit. Los circuitos actuales disipan muchos órdenes de magnitud más por operación; el módulo~17 retoma la cuestión del coste energético del cómputo.

